\section{Simulations}

We now turn to numerical simulations of the model. We pick a standard calibration of 
the investment environment: $\beta = 0.97$ stochastic discount factor, $\alpha = 0.3$ 
capital share, $\delta = 0.04$ depreciation rate, $R = 0.01$ gross interest rate, 
and $\sigma_z = 0.05$ with persistence $\rho_z = 0.9$ for the idiosyncratic 
productivity process (not representative of current conditions but in-line with common
conditions, also not representative of any particular industry). We simulate a population
of $N = 1{,}000$ firms for $T = 300$ periods, with a $2\%$ exit rate per period. Upon
exit, a new firm inherits the capital of the incumbent but starts with a blank slate
of knowledge (i.e., the original prior and no experience). The behavioral distribution is determined to become ergodic by $T = 100$, so we discard 
these first 100 periods as burn-in. We also maintain a "shadow" rational agent to 
benchmark the learning agents' behavior against. Note that because each agent faces
identical environments, they ideally should converge to the same optimal policy, so 
any cross-sectional dispersion in behavior is solely due to differences in learning and 
reasoning dynamics generated by their idiosyncratic experiences.

\subsection{Experience-Only Simulations}

We first simulate using experience-only learning: the firms only update their beliefs 
based on the observed rewards and transitions from their actions, without any reasoning-based deliberation.
In this way we can understand how the learning process alone shapes firm behavior and 
heterogeneity. Second we can study the proper calibration of the reasoning mechanism.

The current hypothesis is as follows: because the belief sequence is a martingale under certainty equivalence, reasoning 
does not shift the conditional mean of $Q^*$. Its sole effect is to compress the 
posterior variance, which in turn tightens the entropy floor governing exploration. 
This compression can be beneficial or detrimental depending on where the firm's 
beliefs stand relative to the truth. When the firm's current policy is close to 
optimal, reasoning usefully stabilizes beliefs and prevents unnecessary 
experimentation from pulling the firm away. When the firm has already deviated 
from the optimal policy, however, the same variance reduction locks in the 
deviation: the tighter posterior forecloses the exploration that would be needed 
to discover and correct the error.


\begin{figure}[ht]
    \centering
    \includegraphics[width=0.95\textwidth]{../figures/learning_dynamics_dashboard.pdf}
    \caption{Learning dynamics for population of firms under experience-only learning 
    ($N=1{,}000$ firms, $T=300$, burn-in $T > 100$). Panel~(a): information gain $\alpha_t$ from experience updates;
    Panel~(b): posterior uncertainty $\Sigma_t$ plotted against the entropy floor constraint $H \cdot \text{Tr}(\Sigma_E)$;
    Panel~(c): Shadow price $\delta_E$ of information; Panel~(d): fraction of firms triggering
    reasoning mechanism plotted against multiple costs $\kappa$. Panel~(e): intensity
    of reasoning measured by dimensions of belief space compressed by reasoning, plotted against $\kappa$. 
    Panel~(f): eigenvalue spectrum vs. reasoning threshold ($\kappa / (H \cdot \delta_E)$) across firms.}
    \label{fig:learning_dynamics_dashboard}
\end{figure}

\ref{fig:learning_dynamics_dashboard} plots the learning dynamics across the population 
of firms under experience-only learning. The first thing to note from Panels~(a) and (b) 
is that the information gain from experience updates stabilizes above zero, and 
posterior uncertainty remains above the entropy floor. The firms suffer from a curse 
of dimensionality in learning. Given they are short-lived with an average lifespan of 
$50$ periods (due to the $2\%$ exit rate), they accumulate roughly $50$ observations of 
the action-value function over a three-dimensional input space $(z, k, i)$. Uniform coverage of a 
$d$-dimensional space requires $n^{1/d}$ observations per dimension, so $50$ 
observations in three dimensions yields approximately $50^{1/3} \approx 3.7$ 
effective observations per axis. Sparse coverage explains why 
the information gain $\alpha_t$ stabilizes above zero and the posterior trace 
remains elevated.

The second row of \ref{fig:learning_dynamics_dashboard} quantifies how 
reasoning would interact with the experience-based learning dynamics. At 
moderate costs ($\kappa_R = 0.0025$), roughly 40\% of firm-periods would 
trigger reasoning, compressing on average 3 eigenvalue dimensions per 
episode. Panel~(f) confirms that the water level $\kappa_R / (H \cdot 
\delta_E)$ sits within the eigenvalue spectrum rather than above or below 
it, so reasoning selectively compresses the dominant directions of 
uncertainty while leaving already-resolved dimensions untouched. As 
reasoning becomes cheaper, the water level drops and more dimensions are 
compressed, accelerating belief convergence but also increasing the risk 
of premature lock-in to policies shaped by early, potentially 
unrepresentative experience.

\begin{table}[ht]
\centering
\small
\begin{tabular}{l ccc ccc c}
\toprule
 & \multicolumn{3}{c}{Investment rate $i/k$} 
 & \multicolumn{3}{c}{Output ratio $y/k$} 
 & Difference \\
\cmidrule(lr){2-4} \cmidrule(lr){5-7} \cmidrule(lr){8-8}
$H$ & Mean & Std & $\overline{\sigma}_{\times}$ 
    & Mean & Std & $\Delta$ 
    & $\Delta\, d/k$ (\%) \\
\midrule
$5 \times 10^{-5}$  & 0.0409 & 0.0444 & 0.0440 & 0.2291 & 0.0887 & 0.0209 & $-11.3$ \\
$1 \times 10^{-4}$  & 0.0411 & 0.0499 & 0.0496 & 0.2364 & 0.1004 & 0.0282 & $-15.3$ \\
$2.5\times 10^{-4}$ & 0.0414 & 0.0531 & 0.0529 & 0.2385 & 0.1060 & 0.0303 & $-16.3$ \\
$5 \times 10^{-4}$  & 0.0413 & 0.0530 & 0.0527 & 0.2488 & 0.1101 & 0.0406 & $-22.6$ \\
$7.5\times 10^{-4}$ & 0.0416 & 0.0540 & 0.0537 & 0.2419 & 0.1098 & 0.0337 & $-18.2$ \\
$1 \times 10^{-3}$  & 0.0415 & 0.0535 & 0.0532 & 0.2399 & 0.1096 & 0.0317 & $-17.0$ \\
$2.5\times 10^{-3}$ & 0.0424 & 0.0607 & 0.0602 & 0.2309 & 0.0979 & 0.0227 & $-10.6$ \\
\midrule
Rational             & 0.0401 & 0.0124 & ---    & 0.2082 & 0.0176 & ---    & $0.0$   \\
\bottomrule
\end{tabular}
\caption{Dispersion statistics under experience-only learning by entropy parameter $H$. 
$N = 1{,}000$ firms, $T = 300$ periods, ergodic sample $t \geq 100$. 
$\text{Std}(i/k)$ is the total standard deviation of the investment rate 
(time and cross-section pooled); $\overline{\sigma}_{\times}(i/k)$ is 
the average cross-sectional standard deviation at each period, isolating 
pure policy dispersion from time-series variation. 
$\Delta d/k$ reports the percentage deviation of mean $d/k$ from the 
rational benchmark.}
\label{tab:dispersion_by_H}
\end{table}

\ref{tab:dispersion_by_H} quantifies how the entropy parameter $H$ shapes the 
ergodic distribution of firm behavior under experience-only learning. Even at the 
lowest $H$, the cross-sectional standard deviation of investment rates 
$\overline{\sigma}_{\times}(i/k)$ is roughly four times the rational benchmark, 
reflecting persistent heterogeneity in learned policies that would not arise under 
rational expectations. As $H$ increases, policy dispersion rises monotonically: 
higher entropy floors force broader exploration, which exposes firms to more diverse 
(and noisier) regions of the state space, producing greater cross-sectional variation 
in converged beliefs. The mean investment rate $\overline{i/k}$ is largely insensitive 
to $H$, remaining within one percentage point of the rational value across all 
calibrations. This dispersion has a compositional consequence 
for the reported ratios. Because some experience-learning firms under-invest and 
shrink, the average capital stock falls below the rational steady state. Under 
concave production ($\alpha < 1$), output per unit capital $y/k = z k^{\alpha - 1}$ 
is decreasing in $k$, so smaller firms mechanically exhibit higher $y/k$. The same 
logic extends to $d/k$: since dividends inherit the concavity of the production 
function through the flow payoff, which is why we higher dividend yield $d/k$ and
negative deviation from the rational benchmark.

\subsection{Experience and Reasoning Simulations}

\begin{figure}[ht]
    \centering
    \includegraphics[width=1.0\textwidth]{../figures/representative_policies_and_envelope.pdf}
    \caption{Expected policy functions ($N=1{,}000$ firms, $T=300$ periods, $2\%$ firm exit rate). 
    Panel~(a): expected policy $\mathbb{E}[k_{t+1} \mid k_t]$ for three 
    representative firms selected from the cross-sectional distribution, 
    with shaded bands showing $\pm 1$ standard deviation of the softmax 
    policy. Panel~(b): cross-sectional policy envelope, showing the mean 
    $\mathbb{E}[k_{t+1} \mid k_t]$ across all firms (solid blue), with 
    the 25th--75th percentile band (dark shading) and 10th--90th percentile 
    band (light shading). The rational benchmark policy is shown in red 
    (dotted). The dash-dotted line indicates the dividend feasibility 
    constraint: the maximum $k_{t+1}$ such that $d_t \geq 0$ at $z = 1$.}
    \label{fig:representative_policies_and_envelope}
\end{figure}

\ref{fig:representative_policies_and_envelope} shows the results of the simulations.
Here we can see that firm heterogeneity endogenously emerges. The population of firms
develops a distribution of policies, with some firms being capital intensive and others 
being more conservative. The cross-sectional policy envelope shows that the mean 
policy is close to the rational benchmark, but there is significant dispersion around 
it. 

\begin{figure}[ht]
    \centering
    \includegraphics[width=0.95\textwidth]{../figures/ergodic_ratios.pdf}
    \caption{Ergodic distributions of firm-level ratios under experience-only learning 
    ($N=1{,}000$ firms, $T=300$, burn-in $T > 100$). Panel~(a): investment rate $i_t/k_t$, 
    Panel~(b): output--capital ratio $y_t/k_t$. Panel~(c): dividend yield $d_t/k_t$.}
    \label{fig:ergodic_ratios}
\end{figure}

The previous figure only shows the resulting distribution of policies at end of simulation,
but to show that the distribution is truly ergodic, we can also look at the 
distribution of firm-level ratios across time and firms. \ref{fig:ergodic_ratios} 
does just this, showing the distribution of the investment rate $i_t/k_t$, 
output--capital ratio $y_t/k_t$, and dividend yield $d_t/k_t$ across all firms and 
periods. Together, these figures show that the model can rationalize investment dispersion
(challenge $\#3$) and lumpy investment dynamics (challenge $\#4$) as emergent properties.

One notable feature is the multi-modality of the ergodic distribution of $i_t/k_t$ 
in panel~(a). This might arise from a shape mismatch between the kernel and the true 
action-value function: the squared-exponential kernel is concave in displacement 
(its covariance saturates at large $\Delta i$), whereas $Q^*$ is convex in 
investment due to quadratic adjustment costs (its curvature increases with 
$|\Delta i|$). As a result, the agent under-generalizes at small deviations from 
its current investment level (believing $Q^*$ varies more than it does locally) 
while over-generalizing at large deviations (believing $Q^*$ varies less than it 
does in the tails).

\begin{figure}[ht]
    \centering
    \includegraphics[width=0.95\textwidth]{../figures/fhp_regression_experience.pdf}
    \caption{Investment regression ($N=1{,}000$ firms, $T=300$, ergodic sample $t > 100$). 
    Panel~(a): naive binned scatter of $i_t/k_t$ on $CF_t/k_t$ without controls. 
    Panel~(b): added-variable plot after partialling out $Q_t$, isolating the 
    coefficient $\hat{\beta}_2$ from the specification 
    $i_t/k_t = a + \beta_1 Q_t + \beta_2 (CF_t/k_t) + \varepsilon_t$. 
    For the rational agent, $Q$ is a sufficient statistic for investment 
    ($\hat{\beta}_1 = 0.42$, $R^2 = 0.99$) and cash flow has no residual 
    explanatory power ($\hat{\beta}_2 \approx 0$).}
    \label{fig:fhp_regression}
\end{figure}

To test for investment dynamics, specifically the presence of sensitivity to $q$ or
cash flow, we run the \cite{FazzariHubbardPetersen1988} regression of $i_t/k_t$ on 
$q_t$ and $CF_t/k_t$. The rational agent, as expected, has a strong sensitivity to 
$q$ and no sensitivity to cash flow after. The learning agent, however, has a 
much weaker sensitivity to $q$ rationalizing the empirical finding that $q$ is a 
weak predictor of investment (challenge $\#1$). However, the learning agent does not
have a significant sensitivity to cash flow after controlling for $q$, which is not
consistent with the empirical finding that cash flow has significant predictive power for
investment (challenge $\#2$). So our model cannot rationalize the cash-flow sensitivity
empirical regularity. These results are shown in \ref{fig:fhp_regression}.

\begin{figure}[ht]
    \centering
    \includegraphics[width=0.95\textwidth]{../figures/lorenz_and_gini_dynamics.pdf}
    \caption{Capital inequality under experience-only learning ($N=1{,}000$ firms, 
    $T=300$, burn-in $T > 100$). Panel~(a): Lorenz curve of time-averaged capital 
    $\bar{k}_j = T_{\text{erg}}^{-1}\sum_{t > t_0} k_{jt}$ across firms, with Gini 
    coefficients reported. Panel~(b): cross-sectional Gini coefficient of $k_t$ at 
    each period.}
    \label{fig:lorenz_gini_dynamics}
\end{figure}

Finally, we examine the distribution of capital across firms. 
Figure~\ref{fig:lorenz_gini_dynamics} plots the Lorenz curve of time-averaged 
capital holdings across firms alongside the cross-sectional Gini coefficient 
at each period. The economy populated by learning and reasoning agents exhibits greater capital 
inequality than the rational reference benchmark, in which identical shadow agents follow 
the optimal policy under the same shock sequences. This heterogeneity arises endogenously from the learning mechanism: 
firms that receive favorable early shock sequences develop locally confident 
(but idiosyncratic) beliefs about optimal investment through both experience 
and reasoning, leading to persistent dispersion in capital accumulation that 
would not emerge under rational expectations.

We now characterize how the reasoning mechanism shapes the learning process. 




